\section{Future Work}

Our experiments focus mainly on the theoretical capabilities of scaling quorum queues horizontally under controlled conditions. In the future, it seems promising to shift the focus towards scenarios that are closer to production environments:

\textbf{Throughput experiment with consumers.} Our linear-capacity experiment measures publish throughput without consumers, using the drop-head overflow strategy to bound memory usage. In a production system, consumers must keep up with the publish rate. Adding consumers introduces additional load on the same single Erlang process that already serializes publishes, which must then also track delivery state, process acknowledgements and manage queue compaction. Measuring throughput with active consumers would provide a more realistic picture of sustainable cluster capacity.

\textbf{Throughput bottleneck} The resource utilization analysis in Section~IV shows that CPU, disk and network are all well below saturation, yet per-node throughput peaks at approximately 34{,}000 msgs/s. It would be interesting to dive deeper into RabbitMQ internals to pinpoint the component that bottlenecks the throughput in our setup.

\textbf{Message variation.} All experiments use a fixed message size of 1\,KB. Message size affects per-message disk overhead, network bandwidth consumption and internal batching efficiency. Smaller messages (e.g. 64 bytes, typical for routing or command payloads) would increase the per-message overhead relative to payload size, while larger messages (e.g. 64\,KB) would shift the bottleneck toward disk and network throughput. Testing across a range of message sizes would reveal whether the observed scaling results hold or change under different payload sizes and characteristics.

\textbf{Fault tolerance under load.} Quorum queues exist primarily to provide data safety during node failures, yet none of our experiments include a failure scenario. Measuring how throughput and latency behave when a node is killed during a benchmark and how quickly the cluster recovers would reveal how different cluster sizes handle failures in practice. In the throughput experiment with quorum size 3, a node failure only affects the queues that had a replica on that node, while the remaining queues continue operating uninterrupted. In the latency experiments with quorum size equal to the cluster size, every node is a voting member and the impact would be more immediate.
